\chapter{Conclusion}

This study showed that PCR-induced chimeric reads in simulated mitochondrial Illumina data can be detected effectively prior to assembly using both engineered-feature machine learning models and sequence-based deep learning models. Across the experiments, all three final models, namely Gradient Boosting, CNN, and BiGRU, were able to distinguish clean from chimeric reads, but they differed in predictive strength, runtime, interpretability, read retention, and downstream assembly behavior.

Among the classical machine learning models, Gradient Boosting emerged as the strongest feature-based model. Its results showed that engineered predictors derived from alignment structure, clipping patterns, and sequence composition are sufficient to support useful chimera detection. Further feature analysis clarified that clipping-based variables, particularly total clipped bases, provided the strongest discriminatory signal, while k-mer discontinuity features also contributed meaningful information. In contrast, microhomology-related predictors and several supplementary-alignment structure summaries contributed relatively little to overall model performance. This was further supported by feature selection and ablation results, which showed that a reduced subset of predictors could perform nearly as well as the full feature set and that removing explicit microhomology descriptors caused only negligible loss in classification quality.

The deep learning models, however, consistently outperformed Gradient Boosting in read-level classification. The CNN demonstrated that direct sequence-based modeling can capture chimera-related patterns beyond those represented by the engineered features alone. The BiGRU produced the strongest read-level classification performance among all evaluated methods, indicating that sequential k-mer representations are especially effective for this task. Its cross-validation performance was also stable across folds.

At the same time, the downstream assembly experiments showed that the best classifier at the read level is not automatically the best filter in practice. Assembly outcomes depended not only on whether chimeric reads were identified correctly, but also on how aggressively reads were removed and how much clean read support remained after filtering. Under lower read depths, read retention became especially important. Gradient Boosting was the most reliable in such cases because it generally preserved more reads while still maintaining contiguous assemblies. The BiGRU showed intermediate behaviour, producing stronger graph simplification than Gradient Boosting but also, in some low-depth and high-chimera rates, reducing read support enough to risk fragmentation. The CNN applied the strongest filtering and therefore sometimes caused fragmentation at low coverage, although under higher-coverage and high-chimera percentage, it produced the best assembly simplification. These results demonstrate that practical utility cannot be based on classification performance alone.

In terms of classifier suitability for pipeline integration, the results indicate a criterion-dependent selection rather than a single universally best model. BiGRU was the most suitable classifier when the priority was read-level predictive performance. Gradient Boosting was the most suitable option when speed, interpretability, and read preservation were prioritized. CNN was useful under selected high-coverage and high-chimera conditions where stronger graph simplification was desired. This supports the modular design of MitoChime, where multiple classifier options can be used depending on the intended analysis context.

Overall, the findings support MitoChime as a useful pre-assembly filtering step for mitochondrial sequencing data. The study showed that model selection should be guided by intended use: a model optimized for read-level classification may differ from one that is more suitable for preserving assembly continuity, improving interpretability, minimizing runtime, or simplifying assembly graphs under specific coverage and chimera conditions.

Future work should extend this evaluation to real mitochondrial and non-mitochondrial sequencing datasets to determine how well the observed patterns generalize under practical laboratory conditions. Additional studies should also examine more realistic class imbalance, broader chimera-generation mechanisms, and sequencing-error profiles that better reflect empirical data. From a modeling perspective, hybrid approaches that combine engineered alignment-derived features with deep sequence encoders may offer a useful direction for improving both robustness and interpretability. Together, these would help determine how well the proposed approaches can support real-world chimera detection workflows beyond the simulated conditions examined in this study.